\section*{Sets and Indices}

\[
P = \{\text{Property\_1},\text{Property\_2},\text{Property\_3},\text{Property\_4}\}
\]

Index:
\[
p \in P
\]

\section*{Parameters}

For each property \(p \in P\):
\begin{itemize}
    \item \(r_p\): annual income of property \(p\).
    \item \(c_p\): purchase cost of property \(p\).
\end{itemize}

Global parameter:
\begin{itemize}
    \item \(B\): total investment budget.
\end{itemize}

From the data:
\[
B = 7{,}000{,}000
\]
\[
(r_{\text{Property\_1}},r_{\text{Property\_2}},r_{\text{Property\_3}},r_{\text{Property\_4}})
=
(12{,}500,35{,}000,23{,}000,100{,}000)
\]
\[
(c_{\text{Property\_1}},c_{\text{Property\_2}},c_{\text{Property\_3}},c_{\text{Property\_4}})
=
(1{,}500{,}000,2{,}100{,}000,2{,}300{,}000,4{,}200{,}000)
\]

The exclusion parameter in the business data indicates that Property\_4 excludes Property\_3; in the implemented model this is represented directly as a pairwise constraint rather than via a separate symbolic parameter.

\section*{Decision Variables}

For each \(p \in P\), define the binary purchase decision
\[
x_p \in \{0,1\}
\qquad \text{(code: \texttt{x[p['name']]})}
\]
where \(x_p = 1\) if property \(p\) is purchased, and \(x_p = 0\) otherwise.

\section*{Objective}

Maximize total annual income:
\[
\max \sum_{p \in P} r_p x_p
\]

Equivalently, with the given data:
\[
\max \; 12500\,x_{\text{Property\_1}}
+ 35000\,x_{\text{Property\_2}}
+ 23000\,x_{\text{Property\_3}}
+ 100000\,x_{\text{Property\_4}}
\]

\section*{Constraints}

Budget limit:
\[
\sum_{p \in P} c_p x_p \le B
\]

With the given data:
\[
1500000\,x_{\text{Property\_1}}
+ 2100000\,x_{\text{Property\_2}}
+ 2300000\,x_{\text{Property\_3}}
+ 4200000\,x_{\text{Property\_4}}
\le 7000000
\]

Mutual exclusion between Property\_4 and Property\_3:
\[
x_{\text{Property\_4}} + x_{\text{Property\_3}} \le 1
\]

Binary restrictions:
\[
x_p \in \{0,1\}
\qquad \forall p \in P
\]

\section*{Notes on Implementation}

The source code constructs a binary linear optimization model in PuLP with one binary variable per property, maximizing total annual income subject to:
\begin{itemize}
    \item one knapsack-style budget constraint, and
    \item one explicit pairwise exclusion constraint, \(\texttt{x['Property\_4']} + \texttt{x['Property\_3']} \le 1\).
\end{itemize}

The model is solved by calling Gurobi through \texttt{GUROBI\_CMD(msg=0)}. No additional valid inequalities, preprocessing reductions, or alternative search logic are implemented in \texttt{current\_heuristic.py}. The exclusion rule from \texttt{general\_parameters.csv} is not modeled as a generic parameterized family of constraints; instead, the code adds this single hard-coded constraint for Property\_4 and Property\_3.
